\section{Hierarchical principals}
\label{sec:hierarchical-principals}

KeyQuorum implements a hierarchical cryptographic principal model in which
organizational ancestry, delegated authority, threshold key reconstruction,
and physical key custody are represented as separate but interacting
security dimensions. The dotted labels seen throughout this manual ---
\code{M}, \code{M.S}, \code{M.S.1} --- are not just a naming convention.
Each one is a \emph{principal}: an entity the system can recognize, place
in an ancestry tree, and reason about independently along each of those
four dimensions. This section makes that model explicit, because it is
what a lot of the behavior in Sections~\ref{sec:devices-and-custody},
\ref{sec:private-bridges}, and \ref{sec:authenticated-updates} is actually
enforcing.

\subsection{What a hierarchical principal is}

In access-control and cryptographic-key-management literature, a
\emph{hierarchical principal} scheme is any design where identities are
arranged in a tree and position in that tree carries meaning: an ancestor
has some form of standing over its descendants that a sibling or an
unrelated node does not. Familiar examples differ sharply in \emph{what}
that standing lets an ancestor do:

\begin{itemize}
  \item In a \textbf{PKI certificate chain}, the hierarchy \emph{is} the
  trust chain, and a parent's compromise is total: whoever holds a CA's
  private key can mint a valid certificate for any descendant, so the
  hierarchy is really one long single point of failure wearing a tree
  shape.
  \item In \textbf{Hierarchical Identity-Based Encryption (HIBE)}, the
  hierarchy is a \emph{key-derivation function}: a parent's master secret
  can compute every descendant's private key directly. Delegation and key
  escrow are the same operation.
  \item In \textbf{Kerberos realms or Unix group trees}, the hierarchy
  governs \emph{grants} (who is allowed to do what) but carries no
  cryptographic material at all.
\end{itemize}

KeyQuorum's dotted-label tree is closer to the third case than the first
two, and that is a deliberate design choice, not an omission. An ancestor
label such as \code{M} or \code{M.S} can \emph{authorize} specific,
narrowly-scoped changes about a descendant --- but it can never
\emph{derive}, \emph{recover}, or \emph{stand in for} that descendant's
private key material. Compromising \code{M}'s signing key lets an attacker
forge authorization letters; it does not, by itself, decrypt anything
\code{M.S.1} holds.

\subsection{Four dimensions, four owners}

The tree of dotted labels is the substrate all four dimensions are indexed
by, but each dimension is a distinct \emph{access structure} --- in the
secret-sharing sense of the term, a rule for which sets of conditions are
authorized --- enforced by its own code and its own invariants:

\begin{enumerate}
  \item \textbf{Organizational ancestry.} A pure naming relation: is
  \code{M.S} an ancestor of, or equal to, \code{M.S.2}? This question is
  answered by string structure alone (dotted-label prefixes), independent
  of what secret, signature, or device sits at either node. It is the one
  piece of the model that every other dimension is built on top of.
  \item \textbf{Delegated authority.} A signature capability: which
  \emph{registered signing keys} may authorize a change addressed to a
  given label? The rule (Section~\ref{sec:authenticated-updates}) is that
  the signer must be the subject itself, or a dotted-label ancestor whose
  signing key this store already has on file. A restructure proposed by a
  non-root label waits for its direct parent's countersignature; a direct
  tree update attempted by that same authorizer instead of a proposal is
  refused outright; and a routine employee reissue by the direct parent
  stays a single signature. Authority here is the power to \emph{authorize
  a documented change}, checked fresh against a signature every time --- it
  is never the power to read or reconstruct what the change is about.
  \item \textbf{Threshold key reconstruction.} A purely mathematical
  property of the Shamir split tree (Section~\ref{sec:quorum-files}): an
  $M$-of-$N$ subset of \emph{sibling} leaf shares reconstructs the secret
  those siblings were split from. This has nothing to do with ancestry or
  authority --- reconstructing \code{M}'s secret from \code{M.S} and
  \code{M.A}'s shares does not require \code{M}'s signature or
  cooperation, and \code{M} cannot reconstruct it alone by virtue of being
  the parent label. \flag{unlock_approval = parent} is an optional,
  additional layer on top of this dimension, not a substitute for it:
  Shamir's threshold must still be met, \emph{and} each contributing leaf's
  direct parent must sign off.
  \item \textbf{Physical custody.} A count of distinct hardware devices
  (Section~\ref{sec:devices-and-custody}): how many separate device ids
  back the shares actually consumed by a reconstruction, independent of
  where those devices' labels sit in the tree. Two slots on one USB drive
  count as one device no matter how far apart their labels are in the
  hierarchy; two key files on two different drives count as two devices
  even if one label is the other's direct parent.
\end{enumerate}

This split is visible in the codebase's own module boundaries, not just in
this manual's organization.

\begin{itemize}
  \item \file{private_bridge.rs} owns the dotted-label ancestry test
  (\code{is_ancestor_or_self}, \code{parent_node_label}).
  \item \file{org_update.rs} and \file{authority.rs} own delegated
  authority.
  \item The Shamir math for threshold reconstruction lives with the split
  tree itself.
  \item \file{keys.rs} and \file{device.rs} own the hardware-key registry
  and physical-device bookkeeping.
\end{itemize}

A change to one dimension should not require touching the others.

\subsection{A worked example}

Take \code{M.S.2}, three labels deep under root \code{M} via \code{M.S}.
The four dimensions give four independent, simultaneous answers about that
one principal:

\begin{center}
\begin{tabular}{@{}lp{0.62\linewidth}@{}}
\toprule
\textbf{Dimension} & \textbf{Answer for \code{M.S.2}} \\
\midrule
Ancestry & \code{M.S} and \code{M} are ancestors; a sibling such as
  \code{M.S.3} is not. \\
Delegated authority & \code{M.S.2} itself, \code{M.S}, or \code{M} may sign
  a reissue naming \code{M.S.2} (Section~\ref{sec:authenticated-updates}),
  provided this store already holds a registered signing key for whichever
  one signs. \\
Threshold reconstruction & \code{M.S.2} is a leaf: it contributes its own
  share toward reconstructing its \emph{parent} \code{M.S}'s secret, an
  $M$-of-$N$ threshold of \code{M.S}'s children's shares --- not a
  signature from \code{M.S} or \code{M}, and not \code{M.S.2}'s share alone
  unless \code{M.S}'s own threshold is 1. \\
Physical custody & If \code{minimum_physical_devices} is 2, the shares
  actually consumed toward that threshold must sit on at least two distinct
  device ids --- regardless of which labels those devices happen to hold. \\
\bottomrule
\end{tabular}
\end{center}

None of these four answers can be derived from another. That independence
is the extension KeyQuorum makes to a plain hierarchical-principal model:
instead of one tree carrying one kind of authority (as a CA chain or a HIBE
key-derivation tree does), the same tree of dotted labels is the shared
index for four composed access structures, each satisfied on its own
terms, each owned by a different piece of the system, and each capable of
tightening independently --- raising a threshold, requiring more physical
devices, adding a signing-authority ancestor, or reshaping the label tree
itself --- without weakening the others.

\begin{noteBox}
This is also why extending the model with a new relationship should not
mean widening \code{org_update.rs} (the orchestration layer) itself. A new
authenticated-update primitive belongs in whichever module owns the table
it actually reads or writes --- \code{key_tree.rs} for tree shape,
\code{private_bridge.rs} for the dotted-label hierarchy and bridge
rosters, or \code{keys.rs} for the hardware-key registry --- so that the
dimension being extended stays legible, and the other three stay exactly
as strict as they were before.
\end{noteBox}
